\lecture{47}{FedNova и справедливость}

\sourcevideo
    {https://www.youtube.com/playlist?list=PLOzRYVm0a65fhKDS8cYIC7YCuVGJ4FOJx}
    {Week 12: Lecture 47: General Update Rule}

Рассматривается федеративная задача
\[
    \min_{x\in\mathbb R^d}F(x),
    \qquad
    F(x)=\sum_{i=1}^{K}p_iF_i(x),
\]
где
\[
    p_i\ge0,
    \qquad
    \sum_{i=1}^{K}p_i=1.
\]
На раунде \(t\) сервер рассылает модель \(x^t\), клиент \(i\)
выполняет \(\tau_i\) локальных шагов, а затем сервер агрегирует
полученные обновления. Общее правило позволяет отделить локальное
направление от масштаба обновления и увидеть, почему FedAvg
перевзвешивает клиентов с разным числом локальных шагов.

\section{Локальная траектория и изменение модели}

Клиент начинает из общей точки
\[
    x_i^{t,0}=x^t
\]
и выполняет рекурсию
\[
    x_i^{t,j+1}
    =
    x_i^{t,j}-\eta g_i^{t,j},
    \qquad
    j=0,\ldots,\tau_i-1,
\]
где \(g_i^{t,j}\) --- локальный стохастический градиент. После
\(\tau_i\) шагов
\[
    x_i^{t,\tau_i}
    =
    x^t
    -
    \eta
    \sum_{j=0}^{\tau_i-1}g_i^{t,j}.
\]
Определим локальное изменение модели
\[
    \Delta_i^t
    =
    x_i^{t,\tau_i}-x^t.
\]
Тогда
\[
    \Delta_i^t
    =
    -\eta
    \sum_{j=0}^{\tau_i-1}g_i^{t,j}.
\]

Обычное федеративное усреднение имеет вид
\[
    x^{t+1}
    =
    \sum_{i=1}^{K}p_i x_i^{t,\tau_i}
    =
    x^t+
    \sum_{i=1}^{K}p_i\Delta_i^t,
\]
то есть
\[
    x^{t+1}
    =
    x^t
    -
    \eta
    \sum_{i=1}^{K}p_i
    \sum_{j=0}^{\tau_i-1}g_i^{t,j}.
\]
Полное изменение клиента растёт не только из-за направления
градиентов, но и из-за числа выполненных локальных шагов.

\section{Нормализованное направление и общее правило}

Пусть локальный метод задаёт неотрицательные коэффициенты
\[
    a_{i,0},\ldots,a_{i,\tau_i-1}
\]
и вектор
\[
    a_i
    =
    (a_{i,0},\ldots,a_{i,\tau_i-1})^{\mathsf T}.
\]
Нормализованное локальное направление определяется формулой
\[
    d_i^t
    =
    \frac1{\lVert a_i\rVert_1}
    \sum_{j=0}^{\tau_i-1}
    a_{i,j}g_i^{t,j},
    \qquad
    \lVert a_i\rVert_1
    =
    \sum_{j=0}^{\tau_i-1}a_{i,j}.
\]
Коэффициенты \(a_{i,j}\) позволяют описывать переменные локальные
шаги и накопление импульса.

Широкий класс федеративных методов записывается как
\[
    x^{t+1}
    =
    x^t
    -
    \eta\tau_{\mathrm{eff}}
    \sum_{i=1}^{K}w_i d_i^t,
\]
где
\[
    w_i\ge0,
    \qquad
    \sum_{i=1}^{K}w_i=1.
\]
Здесь \(d_i^t\) задаёт направление клиента, \(w_i\) --- его
относительный серверный вес, а \(\tau_{\mathrm{eff}}\) --- общий
масштаб раунда. Различные алгоритмы получаются разным выбором
\(a_i\), \(w_i\) и \(\tau_{\mathrm{eff}}\).

Если каждый клиент выполняет один шаг, то
\[
    a_i=(1),
    \qquad
    d_i^t=g_i^{t,0},
    \qquad
    w_i=p_i,
    \qquad
    \tau_{\mathrm{eff}}=1,
\]
и общее правило превращается в обычное агрегирование одного
локального градиента:
\[
    x^{t+1}
    =
    x^t-
    \eta\sum_{i=1}^{K}p_i g_i^{t,0}.
\]

\section{FedAvg и несогласованность цели}

Для FedAvg все внутренние коэффициенты одинаковы:
\[
    a_i=\mathbf1_{\tau_i}.
\]
Тогда
\[
    \lVert a_i\rVert_1=\tau_i,
    \qquad
    d_i^t
    =
    \frac1{\tau_i}
    \sum_{j=0}^{\tau_i-1}g_i^{t,j}.
\]
Это средний локальный градиент клиента. Если положить
\[
    \tau_{\mathrm{eff}}
    =
    \sum_{r=1}^{K}p_r\tau_r
\]
и
\[
    w_i
    =
    \frac{p_i\tau_i}
    {\displaystyle\sum_{r=1}^{K}p_r\tau_r},
\]
то
\[
\begin{aligned}
    \tau_{\mathrm{eff}}
    \sum_{i=1}^{K}w_i d_i^t
    &=
    \sum_{i=1}^{K}
    p_i\tau_i
    \frac1{\tau_i}
    \sum_{j=0}^{\tau_i-1}g_i^{t,j}
    \\
    &=
    \sum_{i=1}^{K}p_i
    \sum_{j=0}^{\tau_i-1}g_i^{t,j}.
\end{aligned}
\]
Следовательно, общее правило точно воспроизводит FedAvg.

Однако исходная целевая функция использует веса \(p_i\), тогда как
эффективные веса в представлении FedAvg равны
\[
    w_i
    =
    \frac{p_i\tau_i}
    {\sum_r p_r\tau_r}.
\]
Поэтому при различных \(\tau_i\) клиент с большим числом локальных
шагов получает больший относительный вес. Это несогласованность между исходной целью и направлением, агрегируемым сервером. При равных
локальных длинах
\[
    \tau_i=\tau
\]
получаем \(w_i=p_i\), и проблема исчезает.

\section{Нормализация FedNova}

FedNova отделяет направление локального обновления от числа
выполненных шагов. Клиент нормализует накопленное изменение:
\[
    \widehat\Delta_i^t
    =
    \frac{\Delta_i^t}{\tau_i}.
\]
Так как
\[
    \Delta_i^t
    =
    -\eta
    \sum_{j=0}^{\tau_i-1}g_i^{t,j},
\]
то
\[
    \widehat\Delta_i^t
    =
    -\eta
    \frac1{\tau_i}
    \sum_{j=0}^{\tau_i-1}g_i^{t,j}.
\]
Клиент передаёт среднее изменение на один локальный шаг, а не полное
накопленное изменение.

Сервер сохраняет исходные веса задачи
\[
    w_i=p_i
\]
и использует масштаб
\[
    \tau_{\mathrm{eff}}
    =
    \sum_{i=1}^{K}p_i\tau_i.
\]
Обновление принимает вид
\[
    x^{t+1}
    =
    x^t
    +
    \tau_{\mathrm{eff}}
    \sum_{i=1}^{K}p_i
    \frac{\Delta_i^t}{\tau_i},
\]
или, через локальные градиенты,
\[
    x^{t+1}
    =
    x^t
    -
    \eta\tau_{\mathrm{eff}}
    \sum_{i=1}^{K}p_i
    \left(
        \frac1{\tau_i}
        \sum_{j=0}^{\tau_i-1}g_i^{t,j}
    \right).
\]
Теперь \(\tau_i\) влияет на общий масштаб раунда, но не изменяет
относительный вес клиента \(p_i\). FedAvg агрегирует
\(\Delta_i^t\), а FedNova --- \(\Delta_i^t/\tau_i\). При
одинаковых \(\tau_i\) оба метода совпадают.

Нормализация устраняет только перевзвешивание по числу локальных
шагов. Она сама по себе не устраняет статистическую неоднородность,
стохастический шум, задержки и различия локальных оптимизаторов.

\section{Меры справедливости}

Минимизация среднего риска может давать хорошее глобальное качество,
но плохой результат для отдельных клиентов. Поэтому дополнительно
рассматривают распределение локальных значений
\[
    F_1(x),\ldots,F_K(x).
\]
Положим
\[
    \overline F(x)
    =
    \frac1K
    \sum_{i=1}^{K}F_i(x).
\]
Простейшая мера разброса --- дисперсия локальных потерь:
\[
    V_F(x)
    =
    \frac1K
    \sum_{i=1}^{K}
    \bigl(F_i(x)-\overline F(x)\bigr)^2.
\]
Малое значение \(V_F\) означает близость клиентских потерь, но не
гарантирует, что сами потери малы. Поэтому дисперсию следует
рассматривать вместе со средней потерей.

Если \(F_i(x)\ge0\) и \(\sum_iF_i(x)>0\), определим
\[
    q_i(x)
    =
    \frac{F_i(x)}
    {\displaystyle\sum_{r=1}^{K}F_r(x)}.
\]
Энтропийная мера равна
\[
    H(x)
    =
    -\sum_{i=1}^{K}q_i(x)\log q_i(x).
\]
При равных значениях \(q_i=1/K\) получаем \(H(x)=\log K\), а при
концентрации значений у небольшого числа клиентов энтропия
уменьшается. Эта мера оценивает равномерность, но не различает
равномерно хорошие и равномерно плохие результаты.

Для неотрицательных показателей качества \(u_i(x)\) индекс Джейна
определяется формулой
\[
    J(x)
    =
    \frac{
        \left(\displaystyle\sum_{i=1}^{K}u_i(x)\right)^2
    }{
        K\displaystyle\sum_{i=1}^{K}u_i(x)^2
    }.
\]
При ненулевых показателях
\[
    \frac1K\le J(x)\le1,
\]
причём \(J(x)=1\) тогда и только тогда, когда все показатели равны.
Обычно \(u_i\) выбирают как точность или положительную полезность;
равенство потерь само по себе ещё не означает высокого качества.

\section{Средняя цель и равномерность}

Стандартная федеративная задача минимизирует
\[
    \sum_{i=1}^{K}p_iF_i(x),
\]
тогда как справедливая постановка дополнительно стремится уменьшить
различия между клиентами. Эти цели могут конфликтовать: уменьшение
средней потери не обязано улучшать худшего клиента, а строгая
равномерность может ухудшить среднее качество.

FedNova решает более узкую задачу: возвращает серверные веса к
коэффициентам \(p_i\), заданным исходной целевой функцией, при
различных \(\tau_i\). Это исправляет несогласованность, вызванную
вычислительной неоднородностью, но не превращает средневзвешенную цель
в самостоятельный критерий справедливости. Дисперсия, энтропия и
индекс Джейна используются как дополнительные показатели и должны
интерпретироваться вместе с общим качеством модели.
